\section{Symbolic budget and completion}\label{sec:completion}

Choose positive budget shares
\begin{equation}\label{eq:budget-shares}
 q_{\mathrm{blk}},q_{\mathrm{gauss}},q_{\mathrm{sib}}>0,
 \qquad
 q_{\mathrm{blk}}+q_{\mathrm{gauss}}+q_{\mathrm{sib}}<1.
\end{equation}
Define the normalized major allowance
\begin{equation}\label{eq:M0}
 M_0=\frac{c_{\mathrm{maj}}}{K_\sigma}.
\end{equation}
All interface parameters
\[
 \mu,\alpha,\rho,c_w,r_T,\delta_T,C_2,K_\sigma,
 q_{\mathrm{blk}},q_{\mathrm{gauss}},q_{\mathrm{sib}}
\]
are fixed before the bottom scale.  The remaining choices are made in the
following order.

\begin{enumerate}[label=\textbf{\arabic*.}]
\item Set
\begin{equation}\label{eq:eta-choice}
 \eta=\frac{q_{\mathrm{blk}}M_0}{b}.
\end{equation}
Apply Theorem~\ref{thm:global-partition} at $c=c_F$, obtaining
$C_{\mathrm{tail}}(c_F)$ and a scale threshold.

\item Choose $C\ge\max\{1,1/\sqrt{v_*}\}$ so large that
\begin{equation}\label{eq:gaussian-budget}
 bC_{\mathrm{tail}}(c_F)e^{-c_FC^2/2}
 <q_{\mathrm{gauss}}M_0.
\end{equation}

\item Put
\begin{equation}\label{eq:AC-Dsib}
 A_C=2C+3,
 \qquad
 D_{\mathrm{sib}}=
 \frac{q_{\mathrm{sib}}M_0}{bA_C}.
\end{equation}
The affine factor is the ceiling estimate arising from
$N=\lceil C/\sctrl\rceil$ and the eventual inequality $\sctrl\le1$:
\begin{equation}\label{eq:N-bound}
 (2N+1)\sctrl\le2C+3=A_C.
\end{equation}

\item Choose $G\ge1$ so that
\begin{equation}\label{eq:G-choice}
 \beta_b^G\le D_{\mathrm{sib}}.
\end{equation}
The same $G$-prime reservoir is used for every prime divisor of $b$ and is
fixed before the mass family.

\item Only now choose $k_0$ above one finite maximum of all preceding
thresholds: the common analytic supply; avoidance of $T$; control load and
variance scale; local rigidity, exceptions, and corrected boundary hypotheses;
level-set encoding and the two tails; availability and separation of the common
reservoir; pair-pool and greedy-mass conditions; inverse-square comparison;
period and fibre geometry; Taylor disk and aggregate remainder; and the
remaining finite arithmetic certificates.
\end{enumerate}
Every threshold in the final maximum depends only on parameters fixed earlier,
so the construction is non-circular.

Multiplying the two minor estimates by $\sctrl$ and using the choices above
gives the three separate contributions
\begin{align}
 b\eta&=q_{\mathrm{blk}}M_0,\label{eq:budget-block}\\
 bC_{\mathrm{tail}}(c_F)e^{-c_FC^2/2}
 &<q_{\mathrm{gauss}}M_0,\label{eq:budget-gauss}\\
 b(2N+1)\sctrl D_{\mathrm{sib}}
 &\le q_{\mathrm{sib}}M_0.\label{eq:budget-sibling}
\end{align}
Therefore
\begin{align}
 \sum_{h\in S_m}|F(h)|
 &<\frac{(q_{\mathrm{blk}}+q_{\mathrm{gauss}}+q_{\mathrm{sib}})M_0}{\sctrl}\\
 &<\frac{M_0}{\sctrl}
 =\frac{c_{\mathrm{maj}}}{K_\sigma\sctrl}
 \le\frac{c_{\mathrm{maj}}}{\sE}.
 \label{eq:strict-minor}
\end{align}
By \eqref{eq:main-lower}, the real major sum is at least
$c_{\mathrm{maj}}/\sE$.  Lemma~\ref{lem:spectral-positivity} gives a subset
$A\subseteq E$ of positive weight satisfying the target congruence.  The
no-wrap bound and exact mean normalization convert it into
\[
 \sum_{e\in A}\frac1e=\frac1b.
\]
All denominators are distinct squarefree semiprimes and avoid $T$.  This proves
Theorem~\ref{thm:structural-construction}; Lemmas~\ref{lem:small-b} and
\ref{lem:numerator-reduction} give the sufficiency direction of Theorem
\ref{thm:headline}.

The equal-quarter allocation and the identity $2004=4\cdot501$ used in one
formal certificate are recovered by taking
$q_{\mathrm{blk}}=q_{\mathrm{gauss}}=q_{\mathrm{sib}}=1/4$ and the released
value $K_\sigma=501$.  They are a symmetric safety allocation, not additional
mathematical invariants.
